\documentclass{standalone}
\usepackage{tenkz}
\makeatletter
\ExplSyntaxOn
\file_input:n { tenkz-kernel.code.tex }
\bool_new:N \g__tenkz_test_north_route_target_bool
\bool_new:N \g__tenkz_test_east_route_target_bool
\cs_new_eq:NN
  \__tenkz_test_route_port_target_add:nnnn
  \__tenkz_kernel_route_port_target_add:nnnn
\cs_set_protected:Npn \__tenkz_kernel_route_port_target_add:nnnn #1#2#3#4
  {
    \str_if_eq:nnT {#2} {n}
      {\bool_gset_true:N \g__tenkz_test_north_route_target_bool}
    \str_if_eq:nnT {#2} {e}
      {\bool_gset_true:N \g__tenkz_test_east_route_target_bool}
    \__tenkz_test_route_port_target_add:nnnn {#1}{#2}{#3}{#4}
  }
\AtEndDocument
  {
    \bool_if:NF \g__tenkz_test_north_route_target_bool
      {\errmessage{north~route~missed~its~physical~port}}
    \bool_if:NF \g__tenkz_test_east_route_target_bool
      {\errmessage{east~route~missed~its~virtual~port}}
  }
\ExplSyntaxOff
\makeatother
\begin{document}
\tenkzkernel
\begin{tenkz}[rows={wire}, cols=3, bonds=none]
  \tn[at=(1,2), name=P, ports={90:physical}]{P}
  \tnwire[
    kind=string,
    name=route,
    route={n of P},
    crossing=under
  ]{open}{open}
\end{tenkz}
\begin{tenkz}[rows={wire,wire,wire}, cols=1, bonds=none]
  \tn[at=(2,1), name=Q, ports={0:virtual}]{Q}
  \tnwire[
    kind=string,
    name=east-route,
    route={e of Q},
    crossing=over
  ]{open}{open}
\end{tenkz}
\end{document}
